\scanpage{87}
qui se groupent en deux paires \deleted{d'arêtes}
\deleted{\unsure{de deux}} formées d'arêtes antipodiques de $\widetilde{P}_M$.
Ces deux paires correspondent aux
deux faces de \deleted{$\unsure{F|P}$}
\[
P_M=(S_M,F_M)
\]
qui contiennent $a$ :
\[
\begin{aligned}
&\text{Pour toute }a\in A_M,\ \text{il existe exactement}\\
&\text{\uline{deux} faces } f\in F_M\ \text{telles que } f\supset a.
\end{aligned}\tag*{(35)}
\]

On a de plus, pour mémoire :
\[
\left\{\begin{aligned}
&\text{Soit } e_{\deleted{\unsure{M}}}\in\overrightarrow\pi,\quad i=\{\pm e_{\deleted{M}}\},\quad\text{soit }\deleted{\unsure{F_+}}\\
&\overrightarrow F_{M,e_{\deleted{M}}}^{\deleted{\unsure{+}}}\subset\overrightarrow F_{M,i}\\
&\shortparallel\\[-3pt]
&\{\vec f\in\overrightarrow F_{M,i}\mid\vec\tau(\vec f)=e_{\deleted{M}}\}\\
&\text{alors }F_{M,e}\text{ est une ``orientation'' de }P_{M,i},\\
&\text{i.e. si }\vec f,\vec f'\in\overrightarrow F_{M,e}\text{ sont \uline{adjacentes}}\\
&\text{i.e. }|\vec f|\cap|\vec f'|=a\in\unsure{\mathfrak P|\mathbb P|\mathcal P}_2(S_M),\ \text{alors les}\\
&\text{orientations de }a\text{ induites par }\vec f,\vec f'\text{ sont}\\
&\text{opposées}.
\end{aligned}\right.\tag*{(36)}
\]
% slanted
\margin{\uline{NB} Si on interprète $i$ comme un ens. à 2. éléments \unsure{$(\subset\overrightarrow\pi)$|$(\in\overrightarrow\pi)$|$(\subset\overrightarrow\Pi)$}, on a $\underset{\inserted{\text{can}}}{\operatorname{Or}}(P_{M,i})\simeq i$}

On récupère les $F_{M,i}$ à partir des $A_{M,i}$
(mais non $F_M$ à partir de $A_M$) par
\[
F_{M,i}=\{f\in\mathfrak P_3(S_M)\mid\forall a\in\mathfrak P_2(f),\ \text{on a }a\in A_{M,i}\}.\tag*{(37)}
\]
% archive sheet number: 86
